\section{Experiments}
\label{sec:exp}

\begin{figure*}[t]
    \centering
    \includegraphics[width=\textwidth,trim={0 10pt 0 0pt},clip]{figs/feature_space_aircraft.pdf}
    \caption{t-SNE on five models' outputs together of Aircraft datasets after MTCL training: only our model maintains a similar feature distribution to the original CLIP ones with minor shift, while the rest significantly distort the feature space.}
    \label{fig:tsne}
\end{figure*}



\paragraph{Implementation.} We use CLIP~\cite{radford2021learning} model with image encoder ViT-B/16~\cite{dosovitskiy2020image}. We conduct training with AdamW~\cite{loshchilov2017decoupled} optimizer and use a label smoothing~\cite{muller2019does} technique for a better baseline result. For multi-domain task continual learning, we train 1K iterations for each task, while for class-incremental learning, we follow the same evaluation setting in~\cite{douillard2022dytox}. More implementation details can be found in the supplementary material.

\subsection{Main Properties}
\label{subsec:main_properties}

We ablate our method in feature-space in \cref{tab:ablations}, and different choices for parameter-space regularization in \cref{tab:pablate}. Several interesting characteristics are noted.

\paragraph{Continual learning loss.} In \cref{tab:loss_type}, several types of loss for feature space are tested. The Feature Distance penalizing absolute distances achieves a low accuracy. Distillation loss on probability distribution regularizes relative distance in the feature space. With image-only or text-only distillation, the Transfer, Avg., and Last accuracy all improve. A further boost in performance occurs with both of the distillation losses.


\begin{table}[t]
\centering
\caption{Comparison of different methods on MTIL in Order I. 
% \kai{adding one columns to fill the width}
}
\label{tab:result}
\resizebox{\columnwidth}{!}{%
% {\small
\begin{tabular}{@{}lcc|cc|cc@{}}
\toprule
Method & Transfer & $\Delta$ & Avg. & $\Delta$ & Last & $\Delta$  \\ \midrule
CLIP \texttt{ViT-B/16@224px}  & & & & & &  \\
\quad Zero-shot  & \baseline{69.4} &  \baseline{0.0} & \baseline{65.3} & \baseline{0.0} & \baseline{65.3} & \baseline{0.0} \\
\quad Continual Learning & 44.6 & \red{-24.8} & 55.9 & \red{-9.4} & 77.3 & \gre{+12.0} \\
\quad \quad LwF~\cite{li2017learning} & 56.9 & \red{-12.5} & 64.7 & \red{-0.6} & 74.6 & \gre{+9.0} \\
\quad \quad iCaRL~\cite{rebuffi2017icarl} & 50.4 & \red{-19.0} & 65.7 & \gre{+0.4} & 80.1 & \gre{+14.8} \\
\quad \quad LwF-VR~\cite{ding2022don} & 57.2 & \red{-12.2} & 65.1 & \red{-0.2} & 76.6 & \gre{+11.3}  \\
% \quad \quad WC~\cite{kirkpatrick2017overcoming} & \textbf{63.3} & \red{-6.3} & 70.7 & \red{-5.4} & 81.2 & \gre{+15.9} \\
\quad \quad WiSE-FT~\cite{wortsman2022robust} & 52.3 & \red{-17.1} & 60.7 & \red{-4.6} & 77.7 & \gre{+12.4} \\
% \quad \quad WE & 55.9 & \red{-13.5} & 64.7 & \red{-0.6} & 80.8 & \gre{+15.5} \\ 
\midrule
\quad \quad ZSCL$^*$ (Ours) & 62.2 & \red{-7.2} & 72.6 & \gre{+7.3} & \textbf{84.5} & \gre{\textbf{+19.2}} \\
\quad \quad ZSCL (Ours) & \textbf{68.1} & \red{\textbf{-1.3}} & \textbf{75.4} & \gre{\textbf{+10.1}} & 83.6 & \gre{+18.3} \\
\bottomrule
\end{tabular}%
}
\end{table}



\paragraph{Data source for replay.} \cref{tab:text_source,tab:image_source,tab:replay_num,tab:replay_classes} seek the standard for a good reference dataset. As shown in \cref{tab:image_source}, distillation on images of current tasks achieves a good Transfer score. However, it hinders the learning on new tasks and results in a low Avg. and Last score. Images in a specific domain (\eg, Flowers) are also not good choices. General images in ImageNet and Conceptual Caption (CC) datasets are examples of good reference datasets. These images are easily available by a web crawler~\cite{thomee2016yfcc100m}. Text with more semantics can improve performance (shown in \cref{tab:text_source}). When using sentences from the Conceptual Caption dataset~\cite{sharma2018conceptual}, or even sentences randomly generated from the CLIP vocabulary, there is no ground truth target among the texts for the image from the ImageNet dataset. However, they all achieve an improvement due to more semantics. The reference image dataset does not need to be labeled, matched with the text, or seen by the CLIP model. 

Enough semantics in the reference image dataset boosts the distillation performance. In \cref{tab:replay_num,tab:replay_classes}, a smaller number of images and classes all lead to a degradation in the performance. Fewer classes of images in the reference dataset have a worse impact on the performance compared with the image numbers. To keep a reasonable memory buffer, we randomly sample 100k images from ImageNet for MTIL and use texts from CC dataset. For class-incremental learning, conceptual caption dataset's validation set (28k images) is used to avoid information leakage.

\paragraph{Teacher model.} \cref{tab:ref_model} shows the performance of distillation loss with different teacher models. Unlike conventional continual learning, the teacher model should not be the one trained on the previous task; otherwise, the deviation from the initial CLIP in the previous task may be amplified. In contrast, with the initial CLIP as the teacher model, not only is the zero-shot performance improved but the Mean and Last scores are also boosted, indicating that preserving high-quality feature space facilitates continual learning.

% We also experiment with an ensembled teacher model and find it a better choice for the distillation. With weight ensemble, a better teacher model at each task is also a reason for the improved performance.

\begin{table}[t]
\centering
\caption{Comparison of different methods on MTIL in Order II.}
\label{tab:result_o2}
\resizebox{\columnwidth}{!}{%
\begin{tabular}{@{}lcc|cc|cc@{}}
\toprule
Method & Transfer & $\Delta$ & Avg. & $\Delta$ & Last & $\Delta$  \\ \midrule
CLIP \texttt{ViT-B/16@224px} & & & &  \\
\quad Zero-shot & \baseline{65.4} & \baseline{0.0} & \baseline{65.3} & \baseline{0.0} & \baseline{65.3} & \baseline{0.0} \\
\quad Continual Learning & 46.6 & \red{-18.8} & 56.2 & \red{-9.1} & 67.4 & \gre{+2.1} \\
% \quad \quad EWC &  &  &  &  &  &  \\
\quad \quad LwF~\cite{li2017learning} & 53.2 & \red{-12.2} & 62.2 & \red{-5.2} & 71.9 & \gre{+6.6} \\
\quad \quad iCaRL~\cite{rebuffi2017icarl} & 50.9 & \red{-14.5} & 56.9 & \red{-8.4} & 71.6 & \gre{+6.3} \\
\quad \quad LwF-VR~\cite{ding2022don} & 53.1 & \red{-12.3} & 60.6 & \red{-7.4} & 68.3 & \gre{+0.9} \\
% \quad \quad WC~\cite{kirkpatrick2017overcoming} &57.1 & \red{-10.3} &87.9  & \red{-1.4} & 79.9 & \gre{+12.5} \\
\quad \quad WiSE-FT~\cite{wortsman2022robust} & 51.0 & \red{-14.4} & 61.5 & \red{-5.9} & 72.2 & \gre{+6.9} \\
% \quad \quad WE &51.4 & \red{-16.0} &89.2 & \red{-0.1} &78.1 & \gre{7.0} \\
% \quad \quad LwM$-$ &56.6  & \red{-10.8} & \textbf{89.7} & \gre{+0.4} &82.9 & \gre{+10.7}  \\
% \quad \quad LwM$-$ \textit{w.} WC & \textbf{62.8} & \red{-4.6} &88.3 & \red{-1.0} &83.9 & \gre{+16.5}\\
% \quad \quad LwM$-$ \textit{w.} WiSE &59.8  & \red{-7.6} &89.5 & \gre{+0.2}& 83.9 & \gre{+16.5}\\
\midrule
\quad \quad ZSCL$^*$ & 59.8 & \red{-5.6} & 71.8 & \gre{+6.5} & 83.3 & \gre{+18.0} \\
\quad \quad ZSCL & \textbf{64.2} & \red{\textbf{-1.2}} & \textbf{74.5} & \gre{\textbf{+9.2}} & \textbf{83.4} & \gre{\textbf{+18.1}} \\
\bottomrule
\end{tabular}%
}
\end{table}

\begin{table*}[t]
\centering
\caption{Transfer, Avg., and Last scores (\%) of different continue training methods on MTIL benchmark.}
\label{tab:finalacc}
% \resizebox{\textwidth}{!}{%
{
\footnotesize
% \small
% \scriptsize
\begin{tabular}{y{120}*{11}{x{18}}}
\toprule
Method & \rot{Aircraft} & \rot{Caltech101} & \rot{CIFAR100} & \rot{DTD} & \rot{EuroSAT} & \rot{Flowers} & \rot{Food} & \rot{MNIST} & \rot{OxfordPet} & \rot{Cars} & \rot{SUN397} \\ \midrule

CLIP \texttt{ViT-B/16@224px} &  &  &  &  &  &  &  &  &  &  & \\
\quad Zero-shot & 24.3 & 88.4 & 68.2 & 44.6 & 54.9 & 71.0 & 88.5 & 59.4 & 89.0 & 64.7 & 65.2 \\
% \quad Linear Probe & 59.5 & 94.7 & 96.2 & 83.1 & 79.2 & 97.1 & 98.1 & 92.8 & 99.9 & 93.1 & 86.7 & 78.4 \\
\quad Fine-tune & 62.0 & 95.1 & 89.6 & 79.5 & 98.9 & 97.5 & 92.7 & 99.6 & 94.7 & 89.6 & 81.8 \\ \midrule
\midrule

% Continual Learning \\
\quad \textbf{Transfer} \\
\quad \quad Continual-FT & &  67.1 & 46.0 & 32.1 & 35.6 & 35.0 & 57.7 & 44.1 & 60.8 & 20.5 & 46.6 \\
\quad \quad LwF~\cite{li2017learning} & & 74.5 & 56.9 & 39.1 & 51.1 & 52.6 & 72.8 & 60.6 & 75.1 & 30.3 & 55.9 \\
\quad \quad iCaRL~\cite{rebuffi2017icarl} & & 56.6 & 44.6 & 32.7 & 39.3 & 46.6 & 68.0 & 46.0 & 77.4 & 31.9 & 60.5\\
\quad \quad LwF-VR~\cite{ding2022don} & & 77.1 & 61.0 & 40.5 & 45.3 & 54.4 & 74.6 & 47.9 & 76.7 & 36.3 & 58.6 \\
\quad \quad WiSE-FT~\cite{wortsman2022robust} & & 73.5 & 55.6 & 35.6 & 41.5 & 47.0 & 68.3 & 53.9 & 69.3 & 26.8 & 51.9  \\
% \quad \quad WE & & 74.7 & 58.6 & 37.9 & 46.3 & 51.2 & 72.6 & 56.8 & 71.6 & 33.2 & 56.1 \\
\quad \quad Dist. only & & 80.1 & 62.2 & 40.2 & 39.9 & 58.1 & 80.8 & 53.4 & 74.6 & 38.1 & 61.9 \\ 
\quad \quad ZSCL$^*$ & & 78.3 & 64.0 & 42.9 & 45.2 & 63.5 & 84.2 & 56.1 & 78.9 & 44.1 & 64.3 \\
\quad \quad ZSCL & & \textbf{86.0} & \textbf{67.4} & \textbf{45.4} & \textbf{50.4} & \textbf{69.1} & \textbf{87.6} & \textbf{61.8} & \textbf{86.8} & \textbf{60.1} & \textbf{66.8} \\ \midrule

\quad \textbf{Avg.} \\
\quad \quad Continual-FT & 25.5 & 81.5 & 59.1 & 53.2 & 64.7 & 51.8 & 63.2 & 64.3 & 69.7 & 31.8 & 49.7 \\
\quad \quad LwF~\cite{li2017learning} & 36.3 & 86.9 & 72.0 & 59.0 & 73.7 & 60.0 & 73.6 & 74.8 & 80.0 & 37.3 & 58.1 \\
\quad \quad iCaRL~\cite{rebuffi2017icarl} & 35.5 & 89.2 & 72.2 & 60.6 & 68.8 & 70.0 & 78.2 & 62.3 & 81.8 & 41.2 & 62.5 \\
\quad \quad LwF-VR~\cite{ding2022don} & 29.6 & 87.7 & 74.4 & 59.5 & 72.4 & 63.6 & 77.0 & 66.7 & 81.2 & 43.7 & 60.7 \\
\quad \quad WiSE-FT~\cite{wortsman2022robust} & 26.7 & 86.5 & 64.3 & 57.1 & 65.7 & 58.7 & 71.1 & 70.5 & 75.8 & 36.9 & 54.6 \\
% \quad \quad WE & 34.7 & 87.8 & 67.3 & 59.2 & 71.9 & 63.5 & 76.2 & 72.4 & 77.7 & 42.7 & 58.4 \\
\quad \quad Dist. only & 48.1 & 90.6 & 79.8 & 63.2 & 75.6 & 72.5 & 84.7 & 70.2 & 79.8 & 46.9 & 63.7 \\ 
\quad \quad ZSCL$^*$ & \textbf{50.7} & 90.9 & 79.8 & 63.8 & 76.6 & 77.3 & 87.0 & 71.9 & 83.0 & 52.0 & 65.9 \\
\quad \quad ZSCL & 45.1 & \textbf{92.0} & \textbf{80.1} & \textbf{64.3} & \textbf{79.5} & \textbf{81.6} & \textbf{89.6} & \textbf{75.2} & \textbf{88.9} & \textbf{64.7} & \textbf{68.0} \\ \midrule

\quad \textbf{Last} \\
\quad \quad Continual-FT & 31.0 & 89.3 & 65.8 & 67.3 & 88.9 & 71.1 & 85.6 & 99.6 & 92.9 & 77.3 & 81.1 \\
\quad \quad LwF~\cite{li2017learning} & 26.3 & 87.5 & 71.9 & 66.6 & 79.9 & 66.9 & 83.8 & 99.6 & 92.1 & 66.1 & 80.4 \\
\quad \quad iCaRL~\cite{rebuffi2017icarl} & 35.8 & \textbf{93.0} & 77.0 & 70.2 & 83.3 & 88.5 & 90.4 & 86.7 & 93.2 & 81.2 & \textbf{81.9} \\
\quad \quad LwF-VR~\cite{ding2022don} & 20.5 & 89.8 & 72.3 & 67.6 & 85.5 & 73.8 & 85.7 & 99.6 & 93.1 & 73.3 & 80.9 \\
\quad \quad WiSE-FT~\cite{wortsman2022robust} & 27.2 & 90.8 & 68.0 & 68.9 & 86.9 & 74.0 & 87.6 & 99.6 & 92.6 & 77.8 & 81.3 \\
% \quad \quad WE & 34.7 & 87.8 & 67.3 & 59.2 & 71.9 & 63.5 & 76.2 & 72.4 & 77.7 & 42.7 & 58.4 \\
\quad \quad Dist. only & 43.3 & 91.9 & \textbf{81.3} & \textbf{72.4} & \textbf{95.1} & 90.5 & 90.4 & \textbf{99.7} & 92.5 & 85.1 & 81.8 \\ 
\quad \quad ZSCL$^*$ & \textbf{46.0} & 92.3 & 81.2 & \textbf{72.4} & 93.0 & \textbf{92.1} & 90.8 & 99.6 & 93.3 & \textbf{86.6} & 81.7 \\
\quad \quad ZSCL & 40.6 & 92.2 & \textbf{81.3} & 70.5 & 94.8 & 90.5 & \textbf{91.9} & 98.7 & \textbf{93.9} & 85.3 & 80.2 \\
\bottomrule
\end{tabular}%
}
\end{table*}


\paragraph{Parameter-space regularization.} In~\cref{tab:pablate}, we experiment with three different parameter-space regularizations. 
% For , the weight averaging happens after training on each dataset.
We experiment with two variants of WiSE-FT~\cite{wortsman2022robust}: the weighted average between the current model with the initial CLIP or the one at the previous task. The result shows the latter one is a better choice because keeping weight averaging with initial CLIP loses the newly learned knowledge. 
We experiment with different $\alpha$ choices for WiSE, and the best result is reported. While distillation loss improves the whole performance, the parameter-space regularization further protects the zero-shot transfer ability with a higher Transfer. Among the three parameter-space regularizations, only WE achieve a better Last score.
% Weight Ensemble (WE) achieves a better performance compared to WiSE-FT. 
WC greatly improves the Transfer scores with a lower Last score. A combination of the weight consolidation loss with weight ensemble achieves a better tradeoff between Transfer and Last value. While ZSCL$^*$, a variant without the WC loss, obtains the highest Last score, the ZSCL with WC loss outperforms it with 2.8\% Avg. and 5.9\% Transfer scores.
% Weight consolidation (WC)~\cite{kirkpatrick2017overcoming} is a regularization with L2 loss on the parameters $\|\theta-\theta_0\|_2^2$. 
% In \cite{kirkpatrick2017overcoming}, different weights for the regularization are applied according to the difference in the continual learning tasks. In CLIP continual learning, to preserve the foundational knowledge, we simply set the $\theta_0$ to be the initial CLIP model. 
% In~\cref{tab:pablate}, the WiSE achieves a better Transfer and Last score than direct fine-tuning, especially in preserving the zero-shot performance.
% The experiment shows that weight sampling interval $I$ has little effect on the final performance. 
% In \cref{tab:result}, LwM$-$ denotes training with only the feature-space regularization of our method. 
% Although different parameter-space regularization boosts the final performance, the weighted ensemble achieves the best performance in protecting the zero-shot performance in Before accuracy.

\subsection{Multi-domain Task Incremental Learning}
\label{subsec:mtcl}

\cref{tab:result} displays the performance of different methods on the MTIL benchmark, and \cref{tab:finalacc} presents the detailed Transfer, Avg, and Last metrics on each dataset. Zero-shot denotes the zero-shot prediction performance of the initial CLIP model, and Fine-tune means the direct fine-tuning accuracy on each dataset, which can be seen as an upper-bound where no forgetting phenomenon happens. Continual learning uses cross-entropy loss to learn each dataset sequentially, where there is a significant forgetting issue on both zero-shot predictions (Transfer drops by 24.8\%) and newly learned knowledge (Last drops by 9.4\%). Previous methods improve the Last performance slightly and cannot maintain a high zero-shot prediction performance. Without WC, ZSCL$^*$ achieves the best Last scores, outperforming the previous best one by 4.4\%. Our method ZSCL improves 9.1\% on Transfer accuracy, with only 1.3\% drops compared with the initial CLIP model, and achieves a 10.1\% gain in the Avg. accuracy.

Figure~\cref{fig:tsne} provides a visualization of the feature space on Aircraft of the original CLIP and four methods trained at the end of MTCL (t-SNE conducted only once on all features collected). This shows our method is capable of maintaining the pretrained feature distribution with a small averaged feature distance and thus preserving zero-shot performance.

The result of the MTIL method in Order-II is presented in \cref{tab:result_o2}. Our method surpasses previous methods in another order setting of the MTIL benchmark. A similar conclusion holds from the results of MTIL Order-II compared with MTIL Order-I. Our method ZSCL outperforms others by 9.2\% on the Avg. score and 18.1\% on the Last score with only a 1.2\% performance loss on the Transfer score. This shows our approach can work for different orders of the multi-domain task incremental learning. In addition, compared with Order-I, previous methods achieve a much lower Last score (e.g., for Continual-Learning, Order-I has 77.3\%, while Order-II has 65.3\%). With ZSCL, the Last score is similar (83.6\% compared with 83.4\%). This shows our method is more robust towards different training orders.


\begin{table}[t]
  \begin{center}
    \caption{Comparison of state-of-the-art CL methods on CIFAR100 benchmark in class-incremental setting.}
    \label{tab:cifar100_class_incre}
    \resizebox{\columnwidth}{!}{%
    % \small
    \begin{tabular}{y{55}{c}c|cc|cc} 
    \toprule
         & \multicolumn{2}{c}{{ 10 steps}} & \multicolumn{2}{c}{{20 steps}} & \multicolumn{2}{c}{50 steps}  \\
    \textbf{Methods} & \textbf{Avg} & \textbf{Last} & \textbf{Avg} & \textbf{Last} & \textbf{Avg} & \textbf{Last} \\
    \midrule
    % iCaRL~\cite{rebuffi2017icarl} & 65.27 & 50.74 & 61.20 & 43.74 & 56.08 & 36.62 \\
    UCIR~\cite{hou2019learning} & 58.66 & 43.39 & 58.17 & 40.63 & 56.86 & 37.09 \\
    BiC~\cite{wu2019large} & 68.80 & 53.54 & 66.48 & 47.02 & 62.09 & 41.04 \\
    RPSNet~\cite{rajasegaran2019adaptive}  & 68.60 & 57.05 & -  & -  & -  & -  \\
    % WA    & 69.46 & 53.78 & 67.33 & 47.31 & 64.32 & 42.14 \\
    PODNet~\cite{douillard2020podnet} & 58.03 & 41.05 & 53.97 & 35.02 & 51.19 & 32.99 \\
    % DER (w/o P)       & \textbf{75.36} & 65.22 & 74.09 & 62.48 & 72.41 & 59.08 \\
    DER~\cite{yan2021dynamically} & \underline{74.64} & 64.35 & 73.98 & 62.55 & 72.05 & 59.76 \\
    % DyTox~\cite{douillard2022dytox} & 67.33 & 51.68 & 67.30 & 48.45 & 64.39 & 43.47 \\
    DyTox+~\cite{douillard2022dytox} & 74.10 & 62.34 & 71.62 & 57.43 & 68.90 & 51.09 \\
    \midrule
    CLIP~\cite{radford2021learning} & 74.47 & \underline{65.92} & \underline{75.20} & \underline{65.74} & \underline{75.67} & \underline{65.94} \\
    FT & 65.46 & 53.23 & 59.69 & 43.13 & 39.23 & 18.89 \\
    LwF~\cite{li2017learning} & 65.86 & 48.04 & 60.64 & 40.56 & 47.69 & 32.90   \\
    iCaRL~\cite{rebuffi2017icarl} & 79.35 & 70.97 & 73.32 & 64.55 & 71.28 & 59.07 \\
    LwF-VR~\cite{ding2022don} & 78.81 & 70.75 & 74.54 & 63.54 & 71.02 & 59.45 \\
    ZSCL (Ours) & \textbf{82.15} & \textbf{73.65} & \textbf{80.39} & \textbf{69.58}  & \textbf{79.92} & \textbf{67.36}   \\
    \textbf{Impr} & \gre{\textbf{+7.68}} & \gre{\textbf{+7.73}} & \gre{\textbf{+5.19}} & \gre{\textbf{+3.84}} & \gre{\textbf{+3.95}} & \gre{\textbf{+1.42}} \\
    \bottomrule
    \end{tabular}}
\end{center}
\end{table}

\subsection{Class Incremental Learning}
\label{subsec:cil}

We evaluate our methods on conventional class incremental learning.
% , which shows our method's ability in learning new tasks and preventing previously learned knowledge. 
\cref{tab:cifar100_class_incre,tab:tiny_imagenet_class_incre} display results on CIFAR100 and TinyImageNet, respectively. 
We re-implement some previous methods with a CLIP backbone (after CLIP in the table), while others using a special network design cannot be easily adapted. Although zero-shot CLIP prediction achieves a good result on these benchmarks, continual learning with direct fine-tuning or LwF~\cite{li2017learning} degrades the performance greatly, especially under the setting of a large step number. This demonstrates a severe catastrophic forgetting phenomenon in fine-tuning the CLIP model. Our method consistently improves the performance of the CLIP model on both Avg. and Last scores with a large gap towards previous ones. 
